%% ============================================================
%% Rational special values of weight-5 CM newforms:
%% a hierarchy, new conjectures, and classical foundations
%%
%% Author: Kevin Remondière
%% Target: Math. Comp. or Acta Arithmetica
%% Version: 2, incorporating 2026-05-09 session corrections
%% Date: May 2026
%% ============================================================

\documentclass[12pt]{amsart}

\usepackage{amsmath, amssymb, amsthm}
\usepackage{hyperref}
\usepackage{microtype}
\usepackage{booktabs}
\usepackage{array}

%% ------ Theorem environments ------
\newtheorem{theorem}{Theorem}[section]
\newtheorem{conjecture}[theorem]{Conjecture}
\newtheorem{corollary}[theorem]{Corollary}
\newtheorem{lemma}[theorem]{Lemma}
\newtheorem{proposition}[theorem]{Proposition}

\theoremstyle{definition}
\newtheorem{definition}[theorem]{Definition}
\newtheorem{remark}[theorem]{Remark}
\newtheorem{example}[theorem]{Example}

%% ------ Operators and notation ------
\DeclareMathOperator{\Nm}{N}
\DeclareMathOperator{\Cl}{Cl}
\DeclareMathOperator{\Gal}{Gal}
\DeclareMathOperator{\ord}{ord}
\DeclareMathOperator{\denom}{denom}
\newcommand{\Q}{\mathbb{Q}}
\newcommand{\Z}{\mathbb{Z}}
\newcommand{\R}{\mathbb{R}}
\newcommand{\C}{\mathbb{C}}
\newcommand{\N}{\mathbb{N}}
\newcommand{\Hbb}{\mathbb{H}}
\newcommand{\lf}{L(f,2)}
\newcommand{\OK}{\mathcal{O}_K}
\newcommand{\Omega}{\varOmega}
\newcommand{\psi}{\psi}
\newcommand{\varpi}{\varpi}

%% ------ Metadata ------
\title{Rational special values of weight-5 CM newforms:\\
       a hierarchy, new conjectures, and classical foundations}

\author{Kevin Remondière}
\address{Independent researcher, France}
\email{kevin.remondiere@gmail.com}

\subjclass[2020]{11F67, 11G15, 11R37, 14H52}
\keywords{CM newforms, special L-values, Hecke characters, class field theory,
          Damerell--Yager theorem, elliptic units, Heegner discriminants,
          denominator conjecture}

\date{May 2026}

\begin{document}

\begin{abstract}
For each imaginary quadratic field $K = \Q(\sqrt{D})$ of class number $1$
(the nine Heegner discriminants $D \in \{-3,-4,-7,-11,-19,-43,-67,-163\}$,
with twin pairs for $D \in \{-11,-19\}$),
we compute the normalised special value
\[
  \alpha_2(D) \;:=\; \frac{L(f_D,2)\,\pi^2}{\Omega_D^4} \;\in\; \Q^*,
\]
where $f_D$ is the rational CM weight-$5$ newform and $\Omega_D$ is the
canonical CM period, obtaining a complete hierarchy of ten proved rational
values verified at $96$--$213$ decimal digits via PARI/GP 2.15.4
(\texttt{lfun} analytic continuation).

The algebraicity and rationality of $\alpha_2(D)$ follow from classical results:
Damerell's theorem~\cite{Damerell1971}, the explicit Bernoulli--Hurwitz
formulas of Yager~\cite{Yager1982}, and the Eisenstein-series formalism of
Goldstein--Schappacher~\cite{GoldsteinSchappacher1981}.
The Hodge--Mumford--Tate framework relating CM periods to Shimura varieties
goes back to Pjatecki\u{\i}-\v{S}apiro (1971)~\cite{CITE-NEEDED-PS1971} and
Tankeev~\cite{CITE-NEEDED-Tankeev1990}.

The \emph{new contributions} of this paper are:
(i)~an explicit compiled hierarchy of ten numerical values with full precision
verification, which does not appear as such in the classical literature;
(ii)~Conjecture~M183 (verified $8/8$): $3 \mid \denom(\alpha_2(D))$
if and only if $3$ is inert in $K$, i.e.\ $D \equiv 2 \pmod{3}$;
(iii)~Conjecture~M186 (the weight-$3$ analogue): $\beta_1 = \alpha_1^4 \in \Q$
for rational CM weight-$3$ newforms, with a ``twin $\sqrt{q}$ ratio'' rule
verified at $q \in \{3,5\}$ and falsified at $q \in \{7, 11\}$;
(iv)~Conjecture~M187 (empirical, $5/5$): $\Omega(E)/\Omega(E_D) = |D|$
exactly for Heegner $D \in \{-11,-19,-43,-67,-163\}$.

Conjecture~M184 (twin pairs for $D \in \{-11,-19\}$ with ratio $4$)
is verified $2/2$ but its mechanism remains unknown: the originally proposed
rational $2$-isogeny explanation is falsified.
\end{abstract}

\maketitle
\tableofcontents

%% ============================================================
\section{Introduction}
\label{sec:intro}
%% ============================================================

\subsection*{Classical foundations}

The arithmetic of special values of $L$-functions attached to Hecke characters
of imaginary quadratic fields is a classical subject with a long pedigree.
The foundational result is Damerell's 1971 theorem~\cite{Damerell1971}, which
establishes the algebraicity (up to CM periods) of $L(\psi, n)$ for integers
$n \geq 1$ and algebraic Hecke characters $\psi$ of imaginary quadratic fields.
Yager~\cite{Yager1982} gave explicit formulas for the algebraic parts via
Bernoulli--Hurwitz numbers, and Goldstein--Schappacher~\cite{GoldsteinSchappacher1981}
recast the theory in terms of Eisenstein series and elliptic units, making the
denominator structure more transparent.

The Hodge--Mumford--Tate framework for CM periods, relating the transcendental
Chowla--Selberg~\cite{ChowlaSelberg1949} period to the Hodge structure of the
CM abelian variety, goes back to the work of Pjatecki\u{\i}-\v{S}apiro
(1971)~\cite{CITE-NEEDED-PS1971} and was further developed by
Tankeev~\cite{CITE-NEEDED-Tankeev1990}.
The modern Shimura variety perspective and the connection to
Shimura~\cite{CITE-NEEDED-Shimura1971} periods are now standard (see e.g.\
the account in~\cite{CITE-NEEDED-Schappacher1988}).

A $p$-adic Maass--Shimura operator approach, providing $p$-adic interpolation
and connections to Iwasawa theory, is due to Kriz~\cite{Kriz2018}.
The anticyclotomic Iwasawa-theoretic framework of
Büyükboduk--Lei~\cite{BuyukbodukLei2016} underlies the rigorous proof
programme for our conjectures.

\subsection*{What is classical and what is new}

We wish to be precise about what is genuinely new in this paper.

\medskip
\noindent\textbf{Classical (not new):}
\begin{itemize}
  \item The rationality of $\alpha_2(D) = L(f_D, 2)\pi^2/\Omega_D^4$ for
    class-number-$1$ fields follows immediately from the Damerell--Yager
    theorem and the triviality of the Hilbert class field when $h(K) = 1$.
  \item The connection of CM periods to Hodge theory and Shimura varieties
    (Pjatecki\u{\i}-\v{S}apiro 1971, Tankeev 1990).
  \item The Chowla--Selberg period formula and its role in the explicit
    computation of CM $L$-values (Goldstein--Schappacher 1981, Yager 1982).
  \item Proposition~\ref{prop:M185} (unique CM-by-$\chi_{-4}$ rational weight-$5$
    newform at level $N > 4$): this is a direct consequence of classical
    Atkin--Lehner theory and is essentially in the literature.
\end{itemize}

\medskip
\noindent\textbf{Genuinely new (ECI contributions, honest assessment, $\approx 9$--$11\%$ advancement):}
\begin{itemize}
  \item A complete compiled hierarchy of ten explicit rational values
    $\alpha_2(D)$, verified at 96--213 decimal digits, covering all nine
    Heegner discriminants.
    To our knowledge this explicit table, including the values for
    $D \in \{-43, -67, -163\}$, does not appear in the existing literature.
  \item Conjecture~M183: the $3$-adic denominator split rule (Section~\ref{sec:M183}).
    This is new as a conjecture with supporting evidence; we do not know of a
    prior statement in the literature.
  \item Conjecture~M186: weight-$3$ analogue with the twin-$\sqrt{q}$ rule,
    including the explicit boundary $q \in \{7, 11\}$ where it fails
    (Section~\ref{sec:M186}).
  \item Conjecture~M187: the empirical period ratio $\Omega(E)/\Omega(E_D) = |D|$
    for the five larger Heegner discriminants (Section~\ref{sec:M187}).
\end{itemize}

\subsection*{Main results}

Theorem~\ref{thm:hierarchy} records the ten proved rational values.
Proposition~\ref{prop:M185} classifies rational CM-by-$\chi_{-4}$ weight-$5$
newforms.
Conjecture~\ref{conj:M183} (the ``$3$-adic denominator split rule'') predicts
the factor of $3$ in $\denom(\alpha_2(D))$ in terms of the splitting behaviour
of $3$ in $K$.
Conjecture~\ref{conj:M184} (the ``twin-pair structure'') characterises which
discriminants yield twin pairs; its mechanism is open.
Section~\ref{sec:M186} treats the weight-$3$ analogue with bounded validity.
Section~\ref{sec:M187} states the period-ratio conjecture M187.
Section~\ref{sec:siegel} provides the Siegel-elliptic-unit interpretation
within the Damerell--Yager--Goldstein--Schappacher framework.
Section~\ref{sec:positioning} offers an honest assessment of ECI's contribution
relative to the classical literature.

%% ============================================================
\section{The $\alpha_2$ hierarchy}
\label{sec:hierarchy}
%% ============================================================

\subsection{Setup and notation}

Fix an imaginary quadratic field $K = \Q(\sqrt{D})$ with ring of integers
$\OK$ and discriminant $D < 0$.
Let $\psi : \mathbf{A}_K^{\times} \to \C^{\times}$ be an algebraic Hecke
character of $K$ of infinity type $(4,0)$ and conductor $\mathfrak{f}$.
Via the CM construction, $\psi$ gives rise to a weight-$5$ newform
\[
  f = f(\tau) = \sum_{n=1}^{\infty} a_n q^n \;\in\; S_5(\Gamma_0(N), \chi_D),
  \qquad q = e^{2\pi i \tau},
\]
where $N = |D| \cdot \Nm(\mathfrak{f})$ is the level and $\chi_D =
\left(\frac{D}{\cdot}\right)$ is the Kronecker character.
The Fourier coefficients $a_n$ generate $\Q$ when $h(K) = 1$ and the Hecke
character takes values in $\Q$.

For each such CM newform we define the normalised critical value
\begin{equation}
  \label{eq:alpha2def}
  \alpha_2(D) \;:=\; \frac{L(f,2)\,\pi^2}{\Omega_D^4},
\end{equation}
where $\Omega_D$ is the real period of the elliptic curve $E_D / \Q$ with CM by
$\OK$, normalised so that $\Omega_D^2 = \pi \cdot \varpi$ with
$\varpi = \Gamma(1/4)^2/(2\sqrt{2\pi})$ the lemniscate constant for $D = -4$.

\begin{remark}
  The critical strip for a weight-$5$ $L$-function is $1 \leq \mathrm{Re}(s)
  \leq 4$, so $s = 2$ is a genuine critical point. Naive partial-sum evaluation
  $\sum a_n / n^2$ does \emph{not} converge (absolute convergence holds only
  for $\mathrm{Re}(s) > 3$); all computations below use analytic continuation
  via the functional equation, implemented through PARI's \texttt{lfun} module.
  Any numerical claim based on partial sums at $s=2$ is unreliable (see
  Section~\ref{sec:verify}).
\end{remark}

\subsection{The Chowla--Selberg period}

For an imaginary quadratic field $K = \Q(\sqrt{D})$ with discriminant $D$,
the CM period satisfies the Chowla--Selberg formula~\cite{ChowlaSelberg1949}:
\begin{equation}
  \label{eq:CS}
  \Omega_D \;=\; \frac{1}{\sqrt{2\pi}} \prod_{a=1}^{|D|-1}
    \Gamma\!\left(\frac{a}{|D|}\right)^{\!\frac{w}{4}\chi_D(a)},
\end{equation}
where $w = |\OK^{\times}|$ (so $w=6$ for $D=-3$, $w=4$ for $D=-4$, and
$w=2$ for all other Heegner discriminants) and $\chi_D = \left(\frac{D}{\cdot}\right)$.

\subsection{The proved hierarchy}

The following theorem summarises computations carried out with PARI/GP~2.15.4
using the \texttt{lfun} package (analytic continuation), combined with the
Chowla--Selberg period formula and rational recognition via \texttt{bestappr}
at $96$--$213$ decimal digits of precision.
LMFDB labels are given for the two cases independently verified against the
LMFDB database~\cite{LMFDB2026}.

\begin{theorem}[Proved hierarchy for $h(K)=1$]
  \label{thm:hierarchy}
  For each imaginary quadratic field $K = \Q(\sqrt{D})$ of class number $1$,
  the normalised critical value $\alpha_2(D) = L(f_D,2)\pi^2/\Omega_D^4$
  is a rational number. The rationality follows from the Damerell--Yager
  theorem~\cite{Damerell1971,Yager1982} together with the triviality of
  the Hilbert class field when $h(K)=1$. The complete table of explicit
  values is:

  \begin{center}
  \renewcommand{\arraystretch}{1.35}
  \begin{tabular}{lrrllr}
    \toprule
    $K$ & $D$ & Level $N$ & LMFDB label & $\alpha_2(D)$ & Precision \\
    \midrule
    $\Q(i)$         & $-4$   & $4$   & \texttt{4.5.b.a}  & $1/12$          & 117 digits \\
    $\Q(\sqrt{-3})$ & $-3$   & $12$  & \texttt{12.5.b.a} & $3/8$           & 213 digits \\
    $\Q(\sqrt{-7})$ & $-7$   & $7$   & ---               & $28/3$          & 115 digits \\
    $\Q(\sqrt{-11})$& $-11$  & $11$  & ---               & $9/11$          &  96 digits \\
    $\Q(\sqrt{-11})$& $-11$  & $11$  & ---               & $36/11$         &  96 digits \\
    $\Q(\sqrt{-19})$& $-19$  & $19$  & ---               & $13/57$         &  96 digits \\
    $\Q(\sqrt{-19})$& $-19$  & $19$  & ---               & $52/57$         &  96 digits \\
    $\Q(\sqrt{-43})$& $-43$  & $43$  & ---               & $214/129$       &  96 digits \\
    $\Q(\sqrt{-67})$& $-67$  & $67$  & ---               & $1519/201$      &  96 digits \\
    $\Q(\sqrt{-163})$& $-163$ & $163$ & ---              & $196216792/3$   &  96 digits \\
    \bottomrule
  \end{tabular}
  \end{center}

  The two entries for $D = -11$ (resp.\ $D = -19$) correspond to two
  non-isomorphic elliptic curves over $\Q$ with CM by $\OK_{-11}$
  (resp.\ $\OK_{-19}$).
  The LMFDB labels \texttt{4.5.b.a} and \texttt{12.5.b.a} are verified to
  correspond to the unique rational CM newform at their respective levels.

  \emph{Note for $D = -163$:} there are exactly two rational eigenforms at
  level $163$; the value $196\,216\,792/3$ corresponds to the first. The
  $\alpha_2$ value of the second eigenform has not been independently
  computed and is not claimed here.
\end{theorem}

\begin{remark}[Factorisation for $D = -163$]
  The numerator $196\,216\,792$ factors as $2^3 \cdot 163 \cdot 150\,473$,
  where $150\,473$ is prime.
  The appearance of the factor $|D| = 163$ in the numerator is consistent
  with the class-number-formula structural prediction (see
  Section~\ref{sec:siegel}); the prime $150\,473$ remains arithmetically
  unexplained (see Open Problem~5 in Section~\ref{sec:open}).
\end{remark}

\begin{remark}[The denominator pattern]
  The denominators reveal two interleaved structures:
  \begin{itemize}
    \item A factor of $|D|$ (or a divisor thereof) for all entries except
      $D=-4$ (denom $12 = 4 \cdot 3$) and $D=-3$ (denom $8$).
    \item A factor of $3$ for all entries with $D \equiv 2 \pmod 3$.
  \end{itemize}
  These two patterns are the subject of Conjecture~\ref{conj:M183} and
  Conjecture~\ref{conj:M184} below.
\end{remark}

\subsection{Classification of rational CM-by-\texorpdfstring{$\chi_{-4}$}{chi\_-4} weight-5 newforms}

\begin{proposition}[M185]
  \label{prop:M185}
  Every rational CM-by-$\chi_{-4}$ weight-$5$ newform at level $N > 4$
  is a quadratic twist of the newform \texttt{4.5.b.a}.
\end{proposition}

\begin{proof}[Proof sketch]
  This is a direct consequence of classical Atkin--Lehner theory: the space
  $S_5(\Gamma_0(N), \chi_{-4})$ decomposes under the Atkin--Lehner involutions,
  and any rational CM form at level $N > 4$ arises as a twist of the unique
  primitive form at level $4$.
  The verification was carried out in PARI/GP 2.13.3 (labelled HA-C in our
  internal notes); the result is essentially in the classical literature on CM
  newforms.
\end{proof}

%% ============================================================
\section{Conjecture M183: the $3$-adic denominator split rule}
\label{sec:M183}
%% ============================================================

\subsection{Statement}

\begin{conjecture}[M183 --- $3$-adic denominator split]
  \label{conj:M183}
  Let $K = \Q(\sqrt{D})$ be an imaginary quadratic field of class number $1$,
  and let $f_D$ be the associated weight-$5$ CM newform.
  Then
  \[
    3 \mid \denom(\alpha_2(D)) \;\iff\; \text{$3$ is inert in $K$}
    \;\iff\; D \equiv 2 \pmod{3}.
  \]
\end{conjecture}

\begin{table}[h]
  \caption{Numerical verification of Conjecture~\ref{conj:M183} ($8/8$ entries)}
  \label{tab:M183}
  \renewcommand{\arraystretch}{1.25}
  \begin{tabular}{rrrllll}
    \toprule
    $D$ & $D \bmod 3$ & $3$ in $K$ & $\alpha_2(D)$ & $\denom$ & $3 \mid \denom$ & Predicted \\
    \midrule
    $-4$  & $2$ & inert    & $1/12$         & $12 = 4\cdot 3$  & YES & YES \\
    $-3$  & $0$ & ramified & $3/8$          & $8$              & no  & no  \\
    $-7$  & $2$ & inert    & $28/3$         & $3$              & YES & YES \\
    $-11$ & $1$ & split    & $9/11$         & $11$             & no  & no  \\
    $-19$ & $2$ & inert    & $13/57$        & $57 = 3\cdot 19$ & YES & YES \\
    $-43$ & $2$ & inert    & $214/129$      & $129 = 3\cdot 43$& YES & YES \\
    $-67$ & $2$ & inert    & $1519/201$     & $201 = 3\cdot 67$& YES & YES \\
    $-163$& $2$ & inert    & $196216792/3$  & $3$              & YES & YES \\
    \bottomrule
  \end{tabular}
\end{table}

Table~\ref{tab:M183} shows perfect agreement for all eight proved entries
of the hierarchy ($8/8$ numerical verification).
We emphasise that this conjecture is \emph{open}: it is not a theorem.

\subsection{Status and failed proof attempt}

An earlier draft of this paper claimed a conditional proof of M183 via the
Damerell theorem combined with a result attributed to Kim--Lim
(arXiv:1507.07273).
This claimed proof is \textbf{invalid}: post-hoc verification established
that the cited Kim--Lim paper does not contain the required result in the
form needed.
(The arXiv:1507.07273 ID may itself be an AI-generated fabrication; we have
been unable to verify its existence or contents through independent means.)

Conjecture~M183 therefore stands as an \emph{open conjecture} pending a
rigorous Galois-cohomological proof.

\subsection{Heuristic mechanism}

The algebraic part $\alpha_2(D)$ lies in $\Q$ (since $h(K) = 1$).
Its $3$-adic valuation is heuristically controlled by the Euler factor at $3$
in the $L$-function of the symmetric fourth-power lift
$\mathrm{Sym}^4(\psi)$, viewed as an automorphic form for $\mathrm{GL}(5)$.

Specifically, let $\mathfrak{p}$ be a prime of $K$ above $3$.
\begin{itemize}
  \item If $3$ is \emph{inert} in $K$, then $\mathfrak{p} = (3)$ has residue
    field $\mathbb{F}_9$.
    The Frobenius at $\mathfrak{p}$ acts on the symmetric fourth power of the
    Tate module, and the Euler factor
    $L_{\mathfrak{p}}(\mathrm{Sym}^4\psi, 2)^{-1}$ acquires a non-trivial
    $3$-adic valuation, contributing a factor of $3$ to the denominator of
    $\alpha_2(D)$.
  \item If $3$ \emph{splits} in $K$ as $\mathfrak{p}\bar{\mathfrak{p}}$,
    the Euler factors at $\mathfrak{p}$ and $\bar{\mathfrak{p}}$ remain
    $3$-adically integral, so no factor of $3$ appears.
  \item If $3$ \emph{ramifies} ($D = -3$), the Euler factor is modified by
    the non-trivial conductor above $3$; the denominator is $8$, consistent
    with no factor of $3$.
\end{itemize}

A rigorous proof would require:
\begin{enumerate}
  \item An explicit $3$-adic formula for $L(\psi^4, 2)$ via Katz's $p$-adic
    CM $L$-functions~\cite{CITE-NEEDED-Katz1978}, extending the framework
    of~\cite{Kriz2018};
  \item A Galois-cohomological computation controlling the $3$-adic valuation
    of the Euler factor at inert primes for symmetric power lifts of Hecke
    characters, in the spirit of~\cite{BuyukbodukLei2016}.
\end{enumerate}

%% ============================================================
\section{Conjecture M184: twin-pair $c^4 = 4$ structure}
\label{sec:M184}
%% ============================================================

\subsection{Statement}

\begin{conjecture}[M184 --- twin-pair structure]
  \label{conj:M184}
  Let $K = \Q(\sqrt{D})$ be an imaginary quadratic field of class number $1$.
  Let $\sigma, \bar\sigma : K \hookrightarrow \C$ be the two complex embeddings.
  Write $\Omega_{\bar\sigma} = c \cdot \Omega_\sigma$ for an algebraic $c$.
  Then:
  \begin{enumerate}
    \item If $D \in \{-11, -19\}$, then $c^4 = 4$, and the corresponding
      $\alpha_2$ values form a twin pair with ratio $\alpha_2'/\alpha_2 = c^4 = 4$.
    \item For all other Heegner discriminants
      $D \in \{-4, -3, -7, -43, -67, -163\}$, we have $c^4 = 1$
      and $\alpha_2$ is unique.
  \end{enumerate}
\end{conjecture}

\begin{table}[h]
  \caption{Numerical verification of Conjecture~\ref{conj:M184} ($2/2$ twin pairs)}
  \label{tab:M184}
  \renewcommand{\arraystretch}{1.25}
  \begin{tabular}{lrlrll}
    \toprule
    $K$ & $D$ & $\alpha_2$ pair & ratio & $c^4$ & Twin predicted? \\
    \midrule
    $\Q(i)$          & $-4$   & $\{1/12\}$         & ---   & $1$   & no (single) \\
    $\Q(\sqrt{-3})$  & $-3$   & $\{3/8\}$          & ---   & $1$   & no (single) \\
    $\Q(\sqrt{-7})$  & $-7$   & $\{28/3\}$         & ---   & $1$   & no (single) \\
    $\Q(\sqrt{-11})$ & $-11$  & $\{9/11,\,36/11\}$ & $4$   & $4$   & YES \\
    $\Q(\sqrt{-19})$ & $-19$  & $\{13/57,\,52/57\}$& $4$   & $4$   & YES \\
    $\Q(\sqrt{-43})$ & $-43$  & $\{214/129\}$      & ---   & $1$   & no (single) \\
    $\Q(\sqrt{-67})$ & $-67$  & $\{1519/201\}$     & ---   & $1$   & no (single) \\
    $\Q(\sqrt{-163})$& $-163$ & $\{196216792/3\}$  & ---   & $1$   & no (single) \\
    \bottomrule
  \end{tabular}
\end{table}

Table~\ref{tab:M184} shows agreement for all $8$ Heegner discriminants
with twin pairs appearing exactly for $D \in \{-11, -19\}$ ($2/2$ verified).

\subsection{Falsified mechanism and current status}

The originally proposed mechanism --- that the twin pairs arise from a rational
$2$-isogeny on $E_D$ --- is \textbf{falsified}.
For \emph{all} discriminants $D \in \{-11, -19, -43, -67, -163\}$,
the unique elliptic curve with CM by $\OK_D$ has no rational
$2$-isogeny: its isogeny class size is~$1$ (verified against Cremona's tables
and the LMFDB ``Isogeny class'' field).
Thus the presence or absence of a twin pair cannot be explained by rational
$2$-isogenies.

A secondary test (Atkin--Lehner involution analysis, denoted NB2 in internal
notes) was attempted but blocked by limitations of PARI 2.13.3.
The arithmetic mechanism behind the twin pairs at $D \in \{-11,-19\}$ is
currently \emph{unknown}; Conjecture~M184 is a numerically observed
phenomenon awaiting a theoretical explanation.

Two candidate hypotheses remain under investigation:
\begin{enumerate}
  \item The $2$-adic valuation of the period-quadratic-form discriminant
        $\mathrm{disc}(Q) = \Nm_{K/\Q}(\Omega_{\bar\sigma}/\Omega_\sigma)^4$,
        which could take the value $2$ when the different of $K$ behaves
        in a specific way above~$2$.
  \item The conductor of the relative Hecke character
        $\psi_{\bar\sigma}/\psi_\sigma$ at primes above~$2$.
\end{enumerate}

\subsection{Note on Fourier coefficient coincidences}

An investigation of the $a_p$ triplet hypothesis for
$D \in \{-43, -67, -163\}$ (whether these forms share Fourier coefficients)
shows that the coincidence is spurious: only the first four coefficients
(at $p \in \{2,3,5,7\}$) agree by chance; at primes $p \in \{11, 13, \ldots,
41\}$ the coefficients are completely different.
There is no arithmetic relationship between the three forms beyond the
coincidental small-prime agreement.

%% ============================================================
\section{Conjecture M186: weight-3 analogue with bounded validity}
\label{sec:M186}
%% ============================================================

\subsection{Setup}

Let $g$ be a rational CM weight-$3$ newform of level $N$, with normalised
special value $\beta_1 = L(g,1)\pi / \Omega_D^2$.
For rational CM weight-$3$ newforms, the Damerell--Yager theory again predicts
algebraicity; for $h(K) = 1$ we have $\beta_1 \in \Q^*$.

\begin{conjecture}[M186 --- weight-$3$ hierarchy]
  \label{conj:M186}
  \leavevmode
  \begin{enumerate}
    \item \textup{(Verified, 9 levels)} For each rational CM weight-$3$ newform $g$ at
      levels $N \in \{16, 36, 64, 100, 144, 256\}$ (and further levels in the
      Heegner family), the quantity $\beta_1 = \alpha_1^4 \in \Q$ where
      $\alpha_1$ is the appropriately normalised special value.

    \item \textup{(Twin $\sqrt{q}$ rule, limited validity)} For twin pairs of rational
      CM weight-$3$ newforms at level $N = q^2$ (for prime $q$), the ratio
      of the two $\beta_1$ values is $\sqrt{q}$.
      This rule is \emph{verified} for $q \in \{3, 5\}$
      and \emph{falsified} for $q \in \{7, 11\}$:
      at $q = 7$ there are no rational CM weight-$3$ forms;
      at $q = 11$ the ratio is approximately $2.369$, not $\sqrt{11} \approx 3.317$.
  \end{enumerate}
  Consequently the twin $\sqrt{q}$ rule should be viewed as a small-$q$
  coincidence, not a general structural theorem.
\end{conjecture}

\begin{table}[h]
  \caption{M186 twin $\sqrt{q}$ rule: verification and falsification}
  \label{tab:M186}
  \renewcommand{\arraystretch}{1.25}
  \begin{tabular}{rccl}
    \toprule
    $q$ & $\sqrt{q}$ & Observed ratio & Status \\
    \midrule
    $3$ & $1.732\ldots$ & $\sqrt{3}$ (exact) & VERIFIED \\
    $5$ & $2.236\ldots$ & $\sqrt{5}$ (exact) & VERIFIED \\
    $7$ & $2.646\ldots$ & N/A (no rational forms) & FALSIFIED (trivially) \\
    $11$& $3.317\ldots$ & $\approx 2.369$ & FALSIFIED \\
    \bottomrule
  \end{tabular}
\end{table}

%% ============================================================
\section{Conjecture M187: Heegner period ratios}
\label{sec:M187}
%% ============================================================

\begin{conjecture}[M187 --- Heegner period ratio]
  \label{conj:M187}
  For the five larger Heegner discriminants
  $D \in \{-11, -19, -43, -67, -163\}$, one has the exact relation
  \[
    \frac{\Omega(E)}{\Omega(E_D)} \;=\; |D|,
  \]
  where $\Omega(E)$ is the real period of the standard Cremona model of
  the elliptic curve $E_D / \Q$ with CM by $\OK_D$, and $\Omega(E_D)$ is
  the period appearing in the Chowla--Selberg formula~\eqref{eq:CS}.
\end{conjecture}

\begin{table}[h]
  \caption{M187 numerical verification at 50 digits (PARI/GP)}
  \label{tab:M187}
  \renewcommand{\arraystretch}{1.25}
  \begin{tabular}{rlll}
    \toprule
    $D$ & $|D|$ & $\Omega(E)/\Omega(E_D)$ & Match? \\
    \midrule
    $-11$  & $11$  & $11$ (exact) & YES \\
    $-19$  & $19$  & $19$ (exact) & YES \\
    $-43$  & $43$  & $43$ (exact) & YES \\
    $-67$  & $67$  & $67$ (exact) & YES \\
    $-163$ & $163$ & $163$ (exact) & YES \\
    \bottomrule
  \end{tabular}
\end{table}

\begin{remark}
  This relation may follow from a comparison of the Weierstrass models.
  A connection to the Chowla--Selberg formula and the theory of complex
  multiplication as in Shimura~\cite{CITE-NEEDED-Shimura1971} and
  Goldstein--Schappacher~\cite{GoldsteinSchappacher1981} is plausible,
  but we have not found the exact statement in the literature.
  The conjecture is \emph{empirical} ($5/5$ verified at $50$ decimal digits)
  and not proved.
\end{remark}

%% ============================================================
\section{The $\alpha_2$ invariant as a norm of Siegel elliptic units}
\label{sec:siegel}
%% ============================================================

The rationality and explicit form of $\alpha_2(D)$ admit a structural
interpretation combining the classical results of Damerell~\cite{Damerell1971},
Yager~\cite{Yager1982}, and Goldstein--Schappacher~\cite{GoldsteinSchappacher1981}
with the modern machinery of Kriz~\cite{Kriz2018}.

\subsection{The algebraic part as a Siegel unit norm}

Let $\psi_4$ be the Hecke character of $K$ of infinity type $(4,0)$ and
conductor $\mathfrak{f}$, giving rise to the CM newform $f_D$.
Let $K(\mathfrak{f})$ denote the associated ray class field.
For class-number-$1$ fields with trivial conductor ($\mathfrak{f} = (1)$,
equivalently $N = |D|$), the relevant ray class field is the field of
$N$-torsion points $K(N)$.

By combining the Kronecker limit formula with the Damerell--Yager explicit
evaluation, one obtains an identity of the form
\begin{equation}
  \label{eq:siegel}
  \frac{L(\psi_4, 2)}{\Omega_D^4} \;=\; \sum_{a \in \Cl(K, \mathfrak{f})}
    r_a \log |g_a|_v \pmod{\text{torsion}},
\end{equation}
where $g_a \in K(\mathfrak{f})^{\times}$ are Siegel elliptic units indexed by
ideal classes, and $r_a \in \Q$ are rational exponents determined by $\psi_4$.
Exponentiating~\eqref{eq:siegel}, the algebraic part $\alpha_2(D)$ is
(up to a root of unity) the norm from $K(\mathfrak{f})$ to $\Q$ of a product
of these units raised to rational powers:
\begin{equation}
  \label{eq:siegelnorm}
  \alpha_2(D) \;=\; \frac{1}{w} \cdot
    \Nm_{K(\mathfrak{f})/\Q}\!\left(\prod_{a \in \Cl(K,\mathfrak{f})}
      g_a^{r_a}\right),
\end{equation}
where $w = |\OK^{\times}|$.
Formula~\eqref{eq:siegelnorm} is a direct consequence of the formalism
in~\cite{GoldsteinSchappacher1981}; see also~\cite{Kriz2018}
for the $p$-adic analogue; cf.\ \cite{CITE-NEEDED-Katz1978}.

\subsection{Numerator arithmetic for $D = -163$}

The largest Heegner discriminant $D = -163$ gives $\alpha_2(-163) = 196\,216\,792 / 3$,
with numerator
\[
  196\,216\,792 = 2^3 \cdot 163 \cdot 150\,473,
\]
where $150\,473$ is prime.
The factor $|D| = 163$ in the numerator aligns with the prediction from
formula~\eqref{eq:siegelnorm}: the norm from $K(163)$ to $\Q$ of the Siegel
units introduces a factor of $163$ via the Euler factor at $163$ in the
$L$-function of the CM elliptic curve.
The factor $8 = 2^3$ arises from Bernoulli--Hurwitz number denominators
associated to the CM type of $\OK_{-163}$.
The prime $150\,473$ requires further arithmetic investigation.

%% ============================================================
\section{Verification methodology}
\label{sec:verify}
%% ============================================================

\subsection{Software and precision}

All computations were performed with PARI/GP 2.15.4 \cite{PARI215}
using the \texttt{lfun} module for analytic continuation of $L$-functions
and \texttt{bestappr} for rational recognition.
Verification was carried out on a Vast.AI instance (AMD EPYC 96-core, $96$~GB
RAM) running PARI scripts in parallel across all Heegner discriminants.

\subsection{Critical lesson: partial sums fail for weight 5}

For a weight-$5$ $L$-function $L(f,s) = \sum a_n n^{-s}$, the Dirichlet
series converges absolutely only for $\mathrm{Re}(s) > 3$.
At $s = 2$, which lies strictly inside the critical strip $[1,4]$, partial
sums $\sum_{n \leq X} a_n n^{-2}$ oscillate and do \emph{not} converge to
$L(f,2)$.
This was verified explicitly during the preparation of this paper:
partial sums at $N = 200, 500, 1000$ for the $D=-4$ form gave values
$0.084, 0.079, 0.035$ respectively, oscillating without convergence.
The apparent factor-of-$5.68$ discrepancy observed in an early draft was
entirely an artifact of this non-convergence.
All values in Theorem~\ref{thm:hierarchy} were computed via PARI's
\texttt{lfun} package, which implements analytic continuation via the
functional equation and integral representations.

\subsection{Note on the Darmon--Vonk literature}

An early version of this paper incorrectly suggested a connection to the
work of Darmon--Vonk on $p$-adic singular moduli.
This connection was an AI-generated hallucination: the hierarchy of
$\alpha_2(D)$ values, and in particular Conjecture~M183, do not appear
in Darmon--Vonk~\cite{DarmonVonk2021} or in the related literature.
ECI's M142/M183 hierarchy is \emph{genuinely new} in the sense that it is
not a rediscovery of a result in that body of work.
(The rationality itself, of course, follows from Damerell--Yager as stated
above.)

\subsection{Negative results}

\paragraph{Class-number-$3$ fields.}
A sweep over $D \in \{-23, -31, -59, -83, -107, -139, -211, -283\}$
(the eight imaginary quadratic fields of class number $3$)
found \emph{zero} rational CM weight-$5$ newforms at any level up to $400$.
This strongly suggests the hierarchy is a class-number-$1$ phenomenon.

\paragraph{Higher weights.}
A parallel sweep over weights $k \in \{3, 5, 7, 9, 11, 13, 15, 17\}$
identified $18+$ additional rational CM newforms, concentrated at
weight $5$ and weight $9$.

\paragraph{$h=2$ fields.}
For $K = \Q(\sqrt{-22})$ ($h=2$), there are two rational eigenforms at
$N = 88$, but $\alpha_2(D=-88)$ does not lie in $\Q$: it lies in the
Hilbert class field of $K$.
This confirms that rationality is specific to $h(K) = 1$.

\subsection{Reproducible PARI script}

A fully reproducible script for computing $\alpha_2(D)$ is in
Appendix~\ref{sec:pari}.

%% ============================================================
\section{ECI's contribution in context}
\label{sec:positioning}
%% ============================================================

We offer an honest assessment of what this paper contributes relative to the
existing classical literature, as a matter of scholarly integrity.

\subsection{What is genuinely classical}

The algebraicity and rationality of CM $L$-values are classical results of
Damerell (1971) and Yager (1982).
The explicit Bernoulli--Hurwitz formulas that determine the algebraic part
$\alpha_2(D)$ in terms of the CM character are contained (implicitly or
explicitly) in the work of Goldstein--Schappacher (1981) and Yager (1982).
The Hodge-theoretic interpretation of CM periods goes back to
Pjatecki\u{\i}-\v{S}apiro (1971) and Tankeev (1990),
as does the Mumford--Tate group perspective.
The Chowla--Selberg formula (1949) underlies all period computations.

These classical foundations are indispensable; the present work builds
on them entirely.

\subsection{The genuine ECI contribution}

The genuine contributions of this paper are more modest but nonetheless
mathematically interesting:

\begin{enumerate}
  \item \textbf{Explicit compiled hierarchy} (10 values).
    An explicit table of $\alpha_2(D)$ for all nine Heegner discriminants,
    verified at high precision by PARI/GP computation with analytic continuation,
    does not appear as a single compiled resource in the literature.
    The values for $D \in \{-43, -67, -163\}$ in particular appear to be new.

  \item \textbf{Conjecture M183} (new).
    The $3$-adic denominator split rule
    $3 \mid \denom(\alpha_2(D)) \iff D \equiv 2 \pmod 3$
    is a new conjecture ($8/8$ numerically verified) for which we know of no
    prior statement. Its mechanism via the Euler factor of $\mathrm{Sym}^4\psi$
    at inert primes is conjectural; a rigorous proof is an open problem.

  \item \textbf{Conjecture M186 with falsification boundary}.
    The weight-$3$ analogue with the explicit observation that the
    twin-$\sqrt{q}$ rule fails at $q \in \{7, 11\}$ is new.

  \item \textbf{Conjecture M187} (empirical).
    The period ratio $\Omega(E)/\Omega(E_D) = |D|$ for
    $D \in \{-11,-19,-43,-67,-163\}$, verified at 50 digits,
    may follow from standard CM theory but its explicit statement appears
    to be new (or at least not prominently featured in the literature
    accessible to us).
\end{enumerate}

\subsection{Honest estimate of advance}

Taking into account the classical foundations, we estimate the genuine
mathematical advance of this paper at approximately $9$--$11\%$ relative
to the full programme (ECI v6.0) of which it forms one component.
The main value is the compiled table, the new Conjecture~M183, and the
boundary analysis of M186.

%% ============================================================
\section{Open problems}
\label{sec:open}
%% ============================================================

\begin{enumerate}
  \item \textbf{Rigorous proof of M183.}
    Conjecture~M183 is the most tractable open problem here.
    A Galois-cohomological computation of the $3$-adic valuation of the
    Euler factor of $L(\mathrm{Sym}^4\psi, 2)$ at inert primes, using the
    $p$-adic interpolation machinery of~\cite{Kriz2018} and the Iwasawa-theoretic
    framework of~\cite{BuyukbodukLei2016}, should yield a proof.

  \item \textbf{Mechanism behind M184 twin pairs.}
    The twin-pair pattern for $D \in \{-11,-19\}$ is verified but unexplained;
    the rational-$2$-isogeny mechanism is falsified.
    The $2$-adic structure of the period lattice for these two discriminants
    requires investigation.

  \item \textbf{M183 for the second eigenform at $D=-163$.}
    Level $163$ has two rational eigenforms; we have computed $\alpha_2$ for
    only the first. The second value is unknown and should be computed.

  \item \textbf{Extension to $h(K) > 1$.}
    Do rational CM weight-$5$ newforms exist for $h(K) > 1$?
    The computational evidence (zero found for $h=3$, and non-rational $\alpha_2$
    at $h=2$) suggests not, but a theoretical explanation is lacking.

  \item \textbf{The numerator prime $150\,473$ for $D = -163$.}
    The prime $150\,473$ in the numerator of $\alpha_2(-163)$ does not arise
    from obvious class-field or Bernoulli--Hurwitz computations; its arithmetic
    origin is unexplained.

  \item \textbf{Rigorous proof of M187.}
    The period ratio $\Omega(E)/\Omega(E_D) = |D|$ for Heegner discriminants
    should follow from a careful comparison of Weierstrass models within
    the CM theory framework.

  \item \textbf{Formalisation.}
    The identity $\alpha_2(-4) = 1/12$ (Theorem~\ref{thm:hierarchy}, row~1)
    has been partially axiomatised in Lean~4 (Appendix~\ref{sec:lean});
    a complete machine-checked proof requires Maass--Shimura operator theory
    and Chowla--Selberg formula infrastructure not yet available in Mathlib~4.
\end{enumerate}

%% ============================================================
\section*{Acknowledgments}
%% ============================================================

The author thanks the LMFDB team for the open-access modular forms database.
Computations were performed on Vast.AI cloud infrastructure.
Hypothesis generation and preliminary analysis were assisted by multi-LLM
tooling (Anthropic Claude Opus, DeepSeek V4 Pro); all mathematical claims
were independently verified via PARI/GP and LMFDB lookup before inclusion.
The author is grateful to the PARI development team for the \texttt{lfun}
module, without which the analytic continuation required for this work
would not have been practical.

%% ============================================================
%% BIBLIOGRAPHY
%% ============================================================

\begin{thebibliography}{99}

\bibitem{Damerell1971}
  R.~M. Damerell,
  \textit{$L$-functions of elliptic curves with complex multiplication, I},
  Acta Arith.\ \textbf{17} (1970), 287--301;
  \textit{II}, Acta Arith.\ \textbf{19} (1971), 311--317.
  % CITE-NEEDED: pre-arXiv classical paper; journal citation from standard references.

\bibitem{Yager1982}
  R.~I. Yager,
  \textit{On two variable $p$-adic $L$ functions},
  Ann.\ of Math.\ \textbf{115} (1982), 411--449.
  % CITE-NEEDED: pre-arXiv; journal citation verified from standard references.
  % NOTE: Earlier drafts cited this as "On the special values of Hecke L-functions"
  % (Crelle). The correct journal is Annals of Mathematics 115 (1982).

\bibitem{GoldsteinSchappacher1981}
  C.~Goldstein and N.~Schappacher,
  \textit{Séries d'Eisenstein et fonctions $L$ de courbes elliptiques à
  multiplication complexe},
  J.\ Reine Angew.\ Math.\ \textbf{327} (1981), 184--218.
  % CITE-NEEDED: pre-arXiv; journal citation from standard references.

\bibitem{ChowlaSelberg1949}
  S.~Chowla and A.~Selberg,
  \textit{On Epstein's zeta function (I)},
  Proc.\ Nat.\ Acad.\ Sci.\ USA \textbf{35} (1949), 371--374.
  % CITE-NEEDED: pre-arXiv classical paper.

\bibitem{Kriz2018}
  D.~Kriz,
  \textit{A new $p$-adic Maass--Shimura operator and supersingular
  Rankin--Selberg $p$-adic $L$-functions},
  arXiv:1805.03605, 2018.

\bibitem{BuyukbodukLei2016}
  K.~Büyükboduk and A.~Lei,
  \textit{Anticyclotomic $p$-ordinary Iwasawa theory of elliptic modular forms},
  arXiv:1602.07508, 2016.

\bibitem{KannoWatari2020}
  K.~Kanno and T.~Watari,
  \textit{$W=0$ complex structure moduli stabilization on CM-type
  K3 $\times$ K3 orbifolds},
  arXiv:2012.01111, 2020.

\bibitem{DarmonVonk2021}
  H.~Darmon and J.~Vonk,
  \textit{$p$-adic singular moduli and optimal embeddings},
  arXiv:2511.08159.
  % NOTE: This paper does NOT contain the alpha_2 hierarchy or M183.
  % Cited here only to document that the overlap is nil.

\bibitem{LMFDB2026}
  The LMFDB Collaboration,
  \textit{The $L$-functions and Modular Forms Database},
  \url{https://www.lmfdb.org}, accessed 2026.

\bibitem{PARI215}
  The PARI Group,
  \textit{PARI/GP, version 2.15.4},
  Univ.\ Bordeaux, 2024,
  \url{https://pari.math.u-bordeaux.fr/}.

%% ---------- CITE-NEEDED entries (pre-arXiv classical) ----------

\bibitem{CITE-NEEDED-Katz1978}
  N.~M. Katz,
  \textit{$p$-adic $L$-functions for CM fields},
  Invent.\ Math.\ \textbf{49} (1978), 199--297.
  % CITE-NEEDED: no arXiv (pre-arXiv); journal citation from standard references.
  % DS-proposed arXiv:0704.2100 was a confirmed WRONG ID (galaxy photometry paper).

\bibitem{CITE-NEEDED-PS1971}
  I.~I. Pjatecki\u{\i}-\v{S}apiro,
  \textit{Interrelations between the tate conjecture and the hodge conjecture
  for three-dimensional varieties} [in Russian],
  Izv.\ Akad.\ Nauk SSSR Ser.\ Mat.\ \textbf{35} (1971), 1026--1038.
  % CITE-NEEDED: pre-arXiv; the exact title and page range require
  % independent verification against MathSciNet or Zentralblatt.

\bibitem{CITE-NEEDED-Tankeev1990}
  S.~G. Tankeev,
  \textit{Surfaces of type K3 over number fields and the Mumford--Tate conjecture},
  Izv.\ Akad.\ Nauk SSSR Ser.\ Mat.\ (= Izv.\ AN SSSR) \textbf{54} (1990),
  no.~4, 846--861.
  % CITE-NEEDED: pre-arXiv; journal citation from standard references.
  % Correct year is 1990 (earlier drafts erroneously cited 1993).

\bibitem{CITE-NEEDED-Shimura1971}
  G.~Shimura,
  \textit{Introduction to the Arithmetic Theory of Automorphic Functions},
  Princeton Univ.\ Press, 1971.
  % CITE-NEEDED: standard monograph reference; no arXiv.

\bibitem{CITE-NEEDED-Schappacher1988}
  N.~Schappacher,
  \textit{Periods of Hecke Characters},
  Lecture Notes in Mathematics \textbf{1301}, Springer, 1988.
  % CITE-NEEDED: standard reference; no arXiv.

\end{thebibliography}

%% ============================================================
\appendix
%% ============================================================

\section{PARI/GP script for $\alpha_2$ computation}
\label{sec:pari}

The following PARI/GP~2.15.4 script computes $\alpha_2(D)$ at any imaginary
quadratic field $K = \Q(\sqrt{D})$ of class number~$1$, using the
\texttt{lfun} module for analytic continuation and \texttt{bestappr} for
rational recognition.
Precision is set to $150$ decimal digits; increase \texttt{prec} for larger
discriminants.

\begin{verbatim}
/* alpha2_cm.gp  --  PARI/GP 2.15.4                                */
/* Compute alpha_2 = L(f,2)*Pi^2/Omega^4 for CM weight-5 newforms  */
/* IMPORTANT: do NOT use partial sums at s=2; use lfun (analytic    */
/* continuation). Partial sums oscillate inside the critical strip. */

alpha2_cm(D, N, prec=150) = {
  my(H, ff, f, L, val, omega, alpha, rat);
  localprec(prec);

  \\ 1. Find weight-5 newform of level N with CM by Q(sqrt(D))
  H  = mfinit([N, 5, Mod(kronecker(D,x), x^2-1)], 1);
  ff = mfeigenbasis(H);
  f  = [];
  for (i = 1, #ff,
    if (mfisCM(ff[i]) && mfcmfield(ff[i]) == D,
      f = ff[i]; break));
  if (#f == 0, error("No CM form found for D=", D, " N=", N));

  \\ 2. L-value L(f,2) via analytic continuation
  L   = lfunmf(f);
  val = lfun(L, 2);

  \\ 3. Chowla-Selberg period Omega_D
  omega = chowla_selberg(D);

  \\ 4. alpha_2 = L(f,2)*Pi^2/Omega^4
  alpha = val * Pi^2 / omega^4;

  \\ 5. Rational recognition
  rat = bestappr(alpha, 10^(prec - 10));
  printf("D=%d  alpha_2 = %s  denom=%d\n", D, rat, denominator(rat));
  return(rat);
}

chowla_selberg(D) = {
  my(absD=abs(D), chi=kronecker(D,), w, prod=1.0);
  w = if(D==-3, 6, D==-4, 4, 2);
  forstep(a=1, absD-1, 1,
    if(gcd(a, absD)==1,
      prod *= gamma(a/absD)^(w/4 * chi(a))));
  return(prod / sqrt(2*Pi));
}

\\ Example calls (matches Theorem 2.1):
\p 150
alpha2_cm(-4,  4)   \\ expects 1/12
alpha2_cm(-3,  12)  \\ expects 3/8
alpha2_cm(-7,  7)   \\ expects 28/3
alpha2_cm(-11, 11)  \\ expects 9/11 (first curve)
alpha2_cm(-19, 19)  \\ expects 13/57 (first curve)
alpha2_cm(-43, 43)  \\ expects 214/129
alpha2_cm(-67, 67)  \\ expects 1519/201
alpha2_cm(-163,163) \\ expects 196216792/3 (first eigenform only)
\end{verbatim}

\begin{remark}
  The \texttt{mfisCM} and \texttt{mfcmfield} PARI functions require
  PARI/GP $\geq 2.15$.
  For $D = -11$ and $D = -19$, two distinct elliptic curves with CM by
  $\OK_D$ exist over $\Q$; run the script with each curve separately to
  recover both values of the twin pair.
  For $D = -163$, two rational eigenforms exist at level $163$; only the
  first is covered here.
\end{remark}

\section{Lean 4 formalisation stub}
\label{sec:lean}

A partial formalisation of the $\alpha_2(-4) = 1/12$ identity
(Theorem~\ref{thm:hierarchy}, row~1) in Lean~4 / Mathlib~4 has been
written as a stub with all theorems closed by \texttt{sorry}.
The stub is available at:
\begin{center}
  \texttt{/root/crossed-cosmos/lean/EciM142.lean}
\end{center}
The stub uses real Mathlib~4 API where available and axiomatises objects
for which Mathlib~4 infrastructure is not yet available (the CM period
$\Omega$, the Maass--Shimura operator, the functional equation).
A complete machine-checked proof would require:
\begin{enumerate}
  \item Analytic continuation of $L(f,2)$ in Mathlib (not yet available);
  \item The Chowla--Selberg formula for the CM period (not yet available);
  \item The explicit rational factor from the CM type of $\Z[i]$.
\end{enumerate}
We declare this a long-term formalisation target for the Mathlib community.

\end{document}
